\documentclass[11pt]{article}
\usepackage[a4paper,margin=26mm]{geometry}
\usepackage{amsmath,amssymb,amsthm,booktabs,hyperref}
\newtheorem{proposition}{Proposition}
\newtheorem{lemma}{Lemma}
\title{Unconditional sphericality and relative-center reduction for $R6_{9,3748}$}
\author{}\date{1 October 2026}
\begin{document}\maketitle

\noindent\textbf{Outcome and scope.}
Every characteristic-zero fusion categorification of the specified ring admits a
spherical structure. Its rank-three subcategory is one of the $\operatorname{PSU}(2)_5$
categories or Galois conjugates; the complementary bimodule is categorically
regular. Consequently the previously conditional fifteen-case relative-center
calculation applies without extra hypotheses. It leaves two one-sided cases,
each of which gives a central lift of the rank-three subcategory.
The companion \texttt{center\_exclusion.tex} uses this central lift to
exclude every categorification. The present note supplies its categorical
reductions. No tensor-factorization theorem for an unbraided ambient
category is asserted.

\paragraph{Input.}
Let $A$ have basis $1,a,b$ with
\[
 a^2=1+b,\qquad ab=a+b,\qquad b^2=1+a+b.
\]
The source ring, including all its structure constants, is
\[
 R=A[X]/(X^2-1-QX),\qquad Q=2+2a+b.
\]
The source labels are $(1,a,b,X,aX,bX)=(0,1,5,2,3,4)$; all simples
are self-dual. The verifier compares the reconstructed tensor with the exact
source tensor. If $\mathcal C$ is a categorification, let $\mathcal D$ be its
subcategory on $1,a,b$, and let $\mathcal M$ be the complementary full
semisimple subcategory. It is a $\mathcal D$-bimodule category.

\begin{lemma}
Every categorification of $A$ has a spherical structure.
\end{lemma}
\begin{proof}
Use the sphericalization $\widetilde{\mathcal D}$ and signed quotient of
\cite[Section 5.1]{OstrikRank3}. Its distinguished central involution
$\varepsilon$ is the negative lift of the unit; the signed quotient imposes
$\varepsilon=-1$. Both lifts of every self-dual simple are self-dual.
Choose lifts of $a,b$. Frobenius reciprocity and the coefficient bounds in
\cite[Section 5.1(b)]{OstrikRank3} give the signed multiplication
\[
 a^2=1+ub,\qquad ab=ua+vb,\qquad b^2=1+va+wb,
 \qquad u,v,w\in\{1,-1\}.
\]
Associativity of $(a^2)b=a(ab)$ forces $uw=1$, hence $w=u$.
Replacing $a$ by $va$ and $b$ by $ub$ makes all three signs positive.
Thus the chosen three lifts form a fusion subcategory mapping equivalently to
$\mathcal D$, which is precisely a spherical section by
\cite[Section 5.1(c)]{OstrikRank3}. No pivotality assumption was made.
Alternatively \cite[Theorem 5.8]{OstrikRank3} applies directly.
\end{proof}

By \cite[Theorem 1.1]{OstrikRank3}, the category $\mathcal D$ is the
even $sl_2$ category at a seventh root of unity, up to the Galois choices.
It admits a nondegenerate braiding, so
\[
 Z(\mathcal D)\simeq\mathcal D\boxtimes\mathcal D^{\mathrm{rev}},
 \qquad K_0(Z(\mathcal D))=A\otimes A.
\]
Here the forgetful map is $u\otimes v\mapsto uv$. This uses a braiding on
$\mathcal D$, not a braiding on $\mathcal C$.

\begin{proposition}
Every categorification of $R$ has a spherical structure.
\end{proposition}
\begin{proof}
Consider the sphericalization of $\mathcal C$. The preceding lemma permits
positive, fusion-closed lifts of $1,a,b$ inside it. Choose a lift $X_+$ and
define the lifts of $aX,bX$ by $a_+X_+,b_+X_+$. These are simple, since their
forgetful images are simple. Associativity and the positive fusion rules of
the lifted $\mathcal D$ then make its entire left action on these three lifts
positive. All twelve simples in the sphericalization are self-dual, so its
Grothendieck ring is commutative and the right action is the same.

In the signed quotient the relation is consequently
\[
 X^2=1+(p+qa+rb)X,
 \qquad p,q\in\{-2,0,2\},\quad r\in\{-1,1\}.
\]
Replacing the three complementary lifts simultaneously by their negative
lifts replaces $(p,q,r)$ by its negative. Thus $r=1$ may be assumed, giving
exactly nine possibilities. All remaining signed products are forced by the
displayed relation and the regular action of $A$; no independent signed
structure constants remain.

Write $S_{p,q}$ for this signed ring and $R$ for the original ring. The
sphericalization tensor is forced to be
\[
 \widetilde N_{(i,e),(j,f)}^{(k,g)}
 =\frac{N_{ij}^k+(-1)^{e+f+g}(S_{p,q})_{ij}^k}{2},
 \qquad e,f,g\in\{0,1\}.
\]
The verifier constructs every rank-twelve tensor, checks all unit, duality,
reciprocity and associativity identities, computes its exact Casimir
polynomial, and factors it over $\mathbb Q$.

If $P_{p,q}(z)$ is the signed rank-six Casimir polynomial, the rank-twelve
Casimir polynomial is
\[
 2^{12}P_{2,2}(z/2)P_{p,q}(z/2).
\]
This equality is checked directly. Thus its formal codegrees include twice
the six roots of $P_{p,q}$. The following table records the coefficient of
$z^5$ and constant coefficient of the monic, irreducible sextic $P_{p,q}$.
Full polynomials and exact nonzero division remainders are in the certificate.
\begin{center}
\begin{tabular}{rrr r l}\toprule
$p$&$q$&$[z^5]P_{p,q}$&$P_{p,q}(0)$&$d$-number test\\\midrule
$-2$&$-2$&$-189$&3296573&fails\\
$-2$&0&$-77$&703493&fails\\
$-2$&2&$-245$&6754013&fails\\
0&$-2$&$-105$&1279733&fails\\
0&0&$-105$&434581&fails\\
0&2&$-385$&2854789&fails\\
2&$-2$&$-133$&1279733&fails\\
2&0&$-245$&3027661&fails\\
2&2&$-637$&5274997&passes\\\bottomrule
\end{tabular}
\end{center}
Each of the first eight rows already violates
$P_{p,q}(0)\mid([z^5]P_{p,q})^6$. Multiplication by the rational number
two preserves the $d$-number property (equivalently, cancels from this
divisibility criterion). Thus each row contradicts
\cite[Theorem 1.2 and Lemma 2.7(v)]{OstrikCodegrees} for the actual
sphericalization. The only possible row is $(p,q,r)=(2,2,1)$.
In that case the six chosen lifts have precisely the positive original
fusion rules, so are closed under tensor products and duality. They give
a fusion subcategory mapping equivalently to $\mathcal C$, hence a spherical
section. This proves the proposition.
\end{proof}

\paragraph{Galois consequence.}
The original formal-codegree polynomial is the irreducible sextic
\[
 P_{2,2}(z)=z^6-637z^5+22295z^4-298116z^3
 +1854944z^2-5274997z+5274997.
\]
For any categorification, the global dimension in its spherical structure is
one of these formal codegrees. Some Galois conjugation sends that root to the
Frobenius--Perron global dimension, giving a pseudo-unitary Galois conjugate.
Consequently existence of a general, pivotal, spherical, or pseudo-unitary
categorification of this particular based ring is equivalent. This does not
say that every categorification is itself pseudo-unitary or unitary.
The standard algebraic-form/Galois-conjugation framework is
\cite{ENO}; the dimension argument is also the one used in
\cite[Lemma 3.4]{OstrikRank3}.

\begin{lemma}
The bimodule $\mathcal M$ is equivalent to the regular
$\mathcal D$-bimodule, and an equivalence can be chosen to send $1$ to $X$.
\end{lemma}
\begin{proof}
The actual functor $L:\mathcal D\to\mathcal M$, $d\mapsto d\otimes X$,
is a left module functor by the ambient associativity constraints.
It maps the three distinct simples bijectively to the three distinct simples
$X,aX,bX$. On simple endomorphism spaces it is the nonzero scalar map,
and all other simple Hom spaces vanish. It is therefore a left-module
equivalence. Equivalently its internal endomorphism algebra at $X$ is $1$.

Transporting the right action through $L$ gives a tensor functor
$\phi:\mathcal D\to\mathcal D$ (with the usual opposite-category convention
for right actions), characterized by
$X\otimes d\simeq L(\phi(d))$. The fusion rules make $\phi$ send each
simple to itself. Hence $\phi$ is a tensor autoequivalence and
$\mathcal M$ is the invertible twisted regular bimodule
${}_{\mathcal D}\mathcal D_{\phi}$.
This proves invertibility categorically, rather than inferring it from the
fusion module alone.

By \cite[Theorem 1.2]{EM},
$\operatorname{BrPic}(\operatorname{Ad}(A_6))=1$, including all its Galois
conjugates. Thus $\mathcal M$ is equivalent to the regular bimodule.
Such an equivalence sends the unit to $X$: if it sent it to $aX$ or $bX$,
left tensoring by $a$ would send a simple to a nonsimple object, impossible
for an equivalence from the regular module.
\end{proof}

\begin{proposition}
The three categorical hypotheses in the previous conditional
relative-center note hold for every categorification of $R$.
\end{proposition}
\begin{proof}
Centers of bimodules are bimodule-functor categories
\cite[Definition 2.1]{GNN}. A regular bimodule equivalence
$E:\mathcal D\to\mathcal M$ induces
$Z(\mathcal D)\simeq Z_{\mathcal D}(\mathcal M)$.
Its unit image is a relative-central object $t=(X,\gamma)$.
With $B=A\otimes A$, the relative-center simple basis is therefore
$B\sqcup Bt$, and its forgetful map sends $t$ to $X$ and
$u\otimes v$ to $uv$.

For completeness, $B$ commutes with $t$ in the Grothendieck ring for a
categorical reason. Use strict notation and the convention
$\gamma_v:vt\to tv$. If $z=(d,\sigma)\in Z(\mathcal D)$, the map
$\gamma_d:dt\to td$ is a morphism in the relative center. The required
identity is
\[
 (\gamma_d\otimes1_v)(1_d\otimes\gamma_v)(\sigma_v\otimes1_t)
 =(1_t\otimes\sigma_v)(\gamma_v\otimes1_d)(1_v\otimes\gamma_d).
\]
It follows by expanding $\gamma_{dv}$ and $\gamma_{vd}$ and applying
naturality of $\gamma$ to $\sigma_v:vd\to dv$. Thus $zt\simeq tz$.
This is also the standard central action of $Z(\mathcal D)$ on the relative
center. The basis element $t$ is self-dual: under the equivalence with
$Z(\mathcal D)$ it is the unique simple forgetting to $X$, since the unit
is the unique simple of $Z(\mathcal D)$ forgetting to $1$.

The nonnegative $B$-part of $t^2$ forgets to $1$, so it is exactly $1$.
The remaining coefficient is $C\in B_{\geq0}$ with $F(C)=Q$. Hence
\[
 K_0(Z_{\mathcal D}(\mathcal C))
 =B[t]/(t^2-1-Ct),\qquad F(C)=2+2a+b.
\]
This is precisely the finite domain of the existing fifteen-case certificate.
\end{proof}

\paragraph{Consequence of the previous arithmetic.}
The already verified $d$-number calculation eliminates thirteen of the
fifteen possibilities. The only remaining possibilities are
\[
 C=Q\otimes1\quad\hbox{or}\quad C=1\otimes Q.
\]
The original arithmetic is unchanged; its applicability is now unconditional.
In either case the corresponding six-simple fusion subcategory of the
relative center maps equivalently to $\mathcal C$. A tensor section
$\mathcal C\to Z_{\mathcal D}(\mathcal C)$ supplies a coherent central
lift $\mathcal D\to Z(\mathcal C)$: for fixed $d\in\mathcal D$, take
its crossing with every $Y\in\mathcal C$ from the relative half-braiding
of the section of $Y$. Naturality in $Y$ and multiplicativity in $Y$ follow
from functoriality and monoidality of the section. Its restriction to
$\mathcal D$ is one of the two nondegenerate braidings.

\paragraph{Role in the complete exclusion.}
The surviving central lift does not supply a braiding on $\mathcal C$.
Instead the companion \texttt{center\_exclusion.tex} applies the
M\"uger factorization theorem inside the braided category $Z(\mathcal C)$.
Its modular centralizer has a dimension multiplicity bound of four, while
a required thirteen-primary twist has a Galois orbit of at least six
objects of the same dimension. This excludes both one-sided relative-center
survivors and, by the spherical/Galois reduction above, every
characteristic-zero categorification. No nongraded extension classification
or associator construction is assumed.

\paragraph{Reproduction.}
Run \texttt{python verify\_sphericality.py --verify}; it replays all nine
rank-twelve tensors and exact polynomial/divisibility calculations.
The script also succeeds with \texttt{python -O}; it has no assertion-based
correctness gates. A generation-and-replay run completed in under two
seconds. The files are confined to this research directory, and the original
package is unchanged.

\begin{thebibliography}{9}
\bibitem{OstrikRank3} V.~Ostrik,
\emph{Pivotal fusion categories of rank 3},
Moscow Math. J. 15 (2015), 373--396,
\href{https://arxiv.org/abs/1309.4822}{arXiv:1309.4822}.
\bibitem{OstrikCodegrees} V.~Ostrik,
\emph{On formal codegrees of fusion categories},
Math. Res. Lett. 16 (2009), 895--901,
\href{https://arxiv.org/abs/0810.3242}{arXiv:0810.3242}.
\bibitem{EM} C.~Edie-Michell,
\emph{The Brauer--Picard groups of the fusion categories coming from the
ADE subfactors}, \href{https://arxiv.org/abs/1709.04721}{arXiv:1709.04721},
Theorem 1.2.
\bibitem{GNN} S.~Gelaki, D.~Naidu, D.~Nikshych,
\emph{Centers of graded fusion categories},
\href{https://arxiv.org/abs/0905.3117}{arXiv:0905.3117}, Section 2.2.
\bibitem{ENO} P.~Etingof, D.~Nikshych, V.~Ostrik,
\emph{On fusion categories},
\href{https://arxiv.org/abs/math/0203060}{arXiv:math/0203060}.
\end{thebibliography}
\end{document}
